\documentclass[11pt]{article}

\usepackage[margin=1in]{geometry}
\usepackage{amsmath,amssymb,mathtools,amsthm,booktabs,microtype,xcolor}
\usepackage{array,enumitem}
\usepackage[colorlinks=true,linkcolor=blue!50!black,urlcolor=blue!50!black]{hyperref}

\title{Kalman Delta Networks 2:\\
Two-Axis Factorized Uncertainty}
\author{}
\date{August 28, 2026}

\newcommand{\R}{\mathbb{R}}
\newcommand{\ve}[1]{\boldsymbol{#1}}
\newcommand{\tr}{\mathsf{T}}
\newcommand{\diag}{\operatorname{diag}}
\newcommand{\Cov}{\operatorname{Cov}}
\newcommand{\KL}{\operatorname{KL}}
\newcommand{\Tr}{\operatorname{Tr}}
\newcommand{\E}{\mathbb{E}}
\newcommand{\Normal}{\mathcal{N}}

\newtheorem{proposition}{Proposition}
\newtheorem{remark}{Remark}

\setlist{itemsep=2pt,topsep=4pt}
\setlength{\emergencystretch}{3em}

\begin{document}
\maketitle

\begin{abstract}
Kalman Delta Networks treat a recurrent key--value memory as the posterior mean of a latent
linear map.  Existing variants retain uncertainty along only one axis.  Diagonal KDN uses a
key-space covariance shared by every value coordinate,
$\ve{I}_{d_v}\otimes\ve{P}_t$, whereas a complementary value-wise isotropic model uses
$\ve{B}_t\otimes\ve{I}_{d_k}$.  This note first derives the value-wise model and its independent
scalar information scans.  It then introduces KDN-2, whose retained covariance is the
two-axis separable family
\[
  \ve{\Sigma}_t=\ve{B}_t\otimes\ve{P}_t,
\]
with both factors positive diagonal.  The exact conditional mean has a simple rank-one form: a
key-wise uncertainty direction multiplies value-wise uncertainty gates.  The exact posterior
covariance, however, is generally not another single Kronecker product.  We derive the exact
posterior, characterize the exceptional closure condition, and obtain the joint reverse-KL
projection back to the separable family.  That projection is a positive matrix-scaling problem,
unique only up to an intrinsic factor gauge.  Consequently, KDN-2 is a principled two-axis
uncertainty model, but it does not inherit the independent M\"obius scans or the single shared
compact-WY memory system of one-axis KDN without an additional approximation.  We also connect
KDN-2 to Gated DeltaNet-2: both expose key- and value-axis modulation, but KDN-2 ties erasure and
writing through one innovation, and their generic algebraic intersection is the scalar KDA rule.
\end{abstract}

\section{Scope and notation}
\label{sec:scope}

Let $K=d_k$ and $V=d_v$.  The recurrent memory
$\ve{S}_t\in\R^{K\times V}$ maps a key to a value.  Its $j$th value column is
$\ve{s}_{t,j}\in\R^K$, and column-major vectorization is
\begin{equation}
  \ve{x}_t=\operatorname{vec}(\widetilde{\ve{S}}_t)
  =\begin{bmatrix}
    \widetilde{\ve{s}}_{t,1}^{\tr}&\cdots&
    \widetilde{\ve{s}}_{t,V}^{\tr}
  \end{bmatrix}^{\tr}.
  \label{eq:vec-memory}
\end{equation}
A tilde denotes the latent random memory, a bare symbol denotes its retained mean, and a hat
denotes the exclusive one-step prediction before the current observation is assimilated.  The
observation model is
\begin{equation}
  \ve{v}_t
  =\widetilde{\ve{S}}_t^{\tr}\ve{k}_t+\ve{e}_t
  =\ve{H}_t\ve{x}_t+\ve{e}_t,
  \qquad
  \ve{H}_t=\ve{I}_V\otimes\ve{k}_t^{\tr},
  \qquad
  \ve{e}_t\sim\Normal(\ve{0},\ve{R}_t).
  \label{eq:observation-model}
\end{equation}
Unless stated otherwise,
$\ve{R}_t=\diag(r_{t,1},\ldots,r_{t,V})\succ0$.  Diagonal observation noise preserves
cross-value independence while allowing different value coordinates to have different
reliability.

The two one-axis uncertainty families are
\begin{equation}
  \underbrace{\ve{I}_V\otimes\ve{P}_t}_{\text{shared key-space uncertainty}}
  \qquad\text{and}\qquad
  \underbrace{\ve{B}_t\otimes\ve{I}_K}_{\text{value-wise isotropic uncertainty}},
  \label{eq:one-axis-families}
\end{equation}
where $\ve{P}_t\in\R^{K\times K}$ and
$\ve{B}_t=\diag(b_{t,1},\ldots,b_{t,V})$.  The first is the family used by Diagonal KDN; the
second is derived next.

\section{Value-wise isotropic uncertainty}
\label{sec:value-isotropic}

\subsection{Retained family and state-space model}

The value-wise isotropic family assigns one scalar key-space variance to every value coordinate:
\begin{equation}
  \widetilde{\ve{s}}_{t,j}\mid\mathcal{F}_t
  \sim\Normal(\ve{s}_{t,j},b_{t,j}\ve{I}_K),
  \qquad
  \Cov\!\left(\operatorname{vec}(\widetilde{\ve{S}}_t)\mid\mathcal{F}_t\right)
  =\ve{B}_t\otimes\ve{I}_K.
  \label{eq:value-family}
\end{equation}
The columns remain conditionally independent, but their uncertainty need not be equal.  Consider
the columnwise process and observation model
\begin{equation}
\begin{aligned}
  \widetilde{\ve{s}}_{t,j}
  &=\ve{D}_t\widetilde{\ve{s}}_{t-1,j}+\ve{w}_{t,j},
  &
  \ve{w}_{t,j}&\sim\Normal(\ve{0},\omega_{t,j}\ve{I}_K),
  \\
  v_{t,j}
  &=\ve{k}_t^{\tr}\widetilde{\ve{s}}_{t,j}+e_{t,j},
  &
  e_{t,j}&\sim\Normal(0,r_{t,j}),
  \qquad
  \ve{D}_t=\diag(\ve{\alpha}_t).
\end{aligned}
  \label{eq:value-state-space}
\end{equation}
The noises are independent across value coordinates and time.  The current shared-noise model is
recovered by tying $\omega_{t,j}=\omega_t$ and $r_{t,j}=r_t$.

The mean prediction is
\begin{equation}
  \widehat{\ve{S}}_t=\ve{D}_t\ve{S}_{t-1}.
  \label{eq:value-mean-prediction}
\end{equation}
The channel-wise transition generally makes the exact predictive covariance anisotropic even
when the retained covariance is isotropic.  Define its average squared transition
\begin{equation}
  a_t=\frac{1}{K}\Tr(\ve{D}_t\ve{D}_t^{\tr})
  =\frac{1}{K}\sum_{i=1}^{K}\alpha_{t,i}^2.
  \label{eq:value-transition-average}
\end{equation}
Projecting the predictive covariance to its average marginal variance gives
\begin{equation}
  \widehat b_{t,j}=a_tb_{t-1,j}+\omega_{t,j},
  \qquad
  \widehat{\ve{P}}_{t,j}=\widehat b_{t,j}\ve{I}_K.
  \label{eq:value-prediction}
\end{equation}
This prediction is exact when
$\ve{D}_t\ve{D}_t^{\tr}=a_t\ve{I}_K$ and is otherwise the covariance-coordinate trace
projection used by Isotropic KDN.

\subsection{Exact conditional mean}

Let
\begin{equation}
  \delta_{t,j}=v_{t,j}-\ve{k}_t^{\tr}\widehat{\ve{s}}_{t,j}.
  \label{eq:value-innovation}
\end{equation}
The predictive cross-covariance and innovation variance are
\begin{equation}
  \Cov(\widetilde{\ve{s}}_{t,j},v_{t,j}\mid\mathcal{F}_{t-1})
  =\widehat b_{t,j}\ve{k}_t,
  \qquad
  \Delta_{t,j}=r_{t,j}+\widehat b_{t,j}\|\ve{k}_t\|_2^2.
  \label{eq:value-joint-moments}
\end{equation}
Therefore column $j$ has scalar gain
\begin{equation}
  \beta_{t,j}=\frac{\widehat b_{t,j}}{\Delta_{t,j}},
  \qquad
  \ve{s}_{t,j}
  =\widehat{\ve{s}}_{t,j}+\beta_{t,j}\ve{k}_t\delta_{t,j}.
  \label{eq:value-gain}
\end{equation}
Every value column writes along the same key direction, but each column uses a different
uncertainty-derived contraction.  In matrix form, with
$\ve{\beta}_t=(\beta_{t,1},\ldots,\beta_{t,V})^{\tr}$,
\begin{equation}
  \ve{S}_t
  =\widehat{\ve{S}}_t
  +\ve{k}_t
   \left[
     \ve{\beta}_t\odot
     \left(\ve{v}_t-\widehat{\ve{S}}_t^{\tr}\ve{k}_t\right)
   \right]^{\tr}.
  \label{eq:value-memory-update}
\end{equation}
Equivalently, grouping the predicted memory columnwise gives
\begin{equation}
  \ve{s}_{t,j}
  =\left(
     \ve{I}_K-\beta_{t,j}\ve{k}_t\ve{k}_t^{\tr}
   \right)\widehat{\ve{s}}_{t,j}
   +\beta_{t,j}\ve{k}_tv_{t,j},
  \qquad j=1,\ldots,V.
  \label{eq:value-memory-column-factorized}
\end{equation}

\subsection{Reverse-KL projection and information scan}

The exact posterior precision for column $j$ is
\begin{equation}
  \ve{C}_{t,j}^{\star}
  =\widehat c_{t,j}\ve{I}_K
   +\frac{1}{r_{t,j}}\ve{k}_t\ve{k}_t^{\tr},
  \qquad
  \widehat c_{t,j}=\widehat b_{t,j}^{-1}.
  \label{eq:value-exact-precision}
\end{equation}
It is dense for a dense key, so it must be projected before the next token.

\begin{proposition}[Value-wise isotropic reverse-KL projection]
Let $\pi_{t,j}$ be the exact Gaussian posterior defined by
Equations~\eqref{eq:value-state-space}--\eqref{eq:value-exact-precision}.  Over
\begin{equation*}
  q_{t,j}=\Normal(\ve{m}_{t,j},b_{t,j}\ve{I}_K),
\end{equation*}
the unique minimizer of $\KL(q_{t,j}\,\|\,\pi_{t,j})$ is
\begin{equation}
  \ve{m}_{t,j}=\ve{s}_{t,j},
  \qquad
  b_{t,j}
  =\frac{K}{\Tr(\ve{C}_{t,j}^{\star})}
  =\left(
    \widehat c_{t,j}
    +\frac{\|\ve{k}_t\|_2^2}{K r_{t,j}}
  \right)^{-1}.
  \label{eq:value-reverse-kl}
\end{equation}
\end{proposition}

\begin{proof}
For a Gaussian target with precision $\ve{C}_{t,j}^{\star}$, the covariance-dependent part of
the reverse KL is
$\tfrac12\{b_{t,j}\Tr(\ve{C}_{t,j}^{\star})-K\log b_{t,j}\}$.  Its derivative vanishes only at
$b_{t,j}=K/\Tr(\ve{C}_{t,j}^{\star})$.  The mean-dependent quadratic is uniquely minimized by
the exact posterior mean in Equation~\eqref{eq:value-gain}.
\end{proof}

Equation~\eqref{eq:value-reverse-kl} is the literal variational update.  As in Isotropic KDN, an
information scale $\mu>0$ can instead be introduced only in the retained post-write precision:
\begin{equation}
  \widehat c_{t,j}
  =\frac{c_{t-1,j}}{a_t+\omega_{t,j}c_{t-1,j}},
  \qquad
  u_{t,j}=\mu\frac{\|\ve{k}_t\|_2^2}{K r_{t,j}},
  \qquad
  c_{t,j}=\widehat c_{t,j}+u_{t,j}.
  \label{eq:value-information-update}
\end{equation}
Here $\mu=1$ is the reverse-KL optimum.  Values $\mu\ne1$ calibrate future protection while
leaving the current gain in Equation~\eqref{eq:value-gain} unchanged.

Substitution gives a scalar M\"obius recurrence for every value coordinate:
\begin{equation}
  c_{t,j}
  =\frac{(1+u_{t,j}\omega_{t,j})c_{t-1,j}+u_{t,j}a_t}
    {\omega_{t,j}c_{t-1,j}+a_t},
  \qquad
  \ve{M}_{t,j}
  =\begin{bmatrix}
    1+u_{t,j}\omega_{t,j} & u_{t,j}a_t\\
    \omega_{t,j} & a_t
  \end{bmatrix}.
  \label{eq:value-mobius}
\end{equation}
Thus the uncertainty stage is $V$ independent associative scalar scans.

\subsection{When the model is genuinely different}

If $b_{0,j}$, $\omega_{t,j}$, and $r_{t,j}$ are tied across $j$, then induction through
Equations~\eqref{eq:value-prediction} and \eqref{eq:value-information-update} keeps every
$b_{t,j}$ and $\beta_{t,j}$ equal.  Realized values cannot break this symmetry because Gaussian
covariance updates do not depend on the realized observation.  At least one value-wise initial
variance, process uncertainty, or observation noise must vary for the model to differ from
Isotropic KDN.

The sequential mean update still costs $O(KV)$, but the transition applied to column $j$ is
\begin{equation}
  \ve{A}_{t,j}
  =(\ve{I}_K-\beta_{t,j}\ve{k}_t\ve{k}_t^{\tr})\ve{D}_t.
  \label{eq:value-column-transition}
\end{equation}
Because $\ve{A}_{t,j}$ varies with $j$, the single compact-WY system shared by all value columns
in current KDN kernels no longer applies.  An exact parallel implementation needs value-batched
systems or a grouped-value approximation.

\section{KDN-2: combined key--value factorized uncertainty}
\label{sec:kdn2}

\subsection{Separable diagonal family}

KDN-2 combines the two one-axis families by retaining
\begin{equation}
\begin{aligned}
  \ve{P}_t&=\diag(p_{t,1},\ldots,p_{t,K})\succ0,
  &
  \ve{B}_t&=\diag(b_{t,1},\ldots,b_{t,V})\succ0,
  \\
  \ve{\Sigma}_t
  &=\ve{B}_t\otimes\ve{P}_t
   =(\ve{B}_t\otimes\ve{I}_K)
    (\ve{I}_V\otimes\ve{P}_t),
  &
  \ve{J}_t
  &=\ve{\Sigma}_t^{-1}
   =\ve{G}_t\otimes\ve{C}_t,
\end{aligned}
  \label{eq:kdn2-family}
\end{equation}
where $\ve{G}_t=\ve{B}_t^{-1}$ and $\ve{C}_t=\ve{P}_t^{-1}$.  Hence the covariance and
information forms are the same family: inversion simply inverts both diagonal factors.  Entry
$(i,j)$ of the latent memory has variance
\begin{equation}
  \operatorname{Var}(\widetilde S_{t,ij}\mid\mathcal{F}_t)=p_{t,i}b_{t,j}.
  \label{eq:kdn2-entry-variance}
\end{equation}
The model has $K+V-1$ covariance degrees of freedom, not $KV$, because
\begin{equation}
  \ve{B}_t\otimes\ve{P}_t
  =(\lambda\ve{B}_t)\otimes(\ve{P}_t/\lambda)
  \qquad\text{for every }\lambda>0.
  \label{eq:kdn2-factor-gauge}
\end{equation}
We fix this internal factor gauge by imposing
\begin{equation}
  \frac{1}{V}\sum_{j=1}^{V}\log b_{t,j}=0
  \qquad\Longleftrightarrow\qquad
  \det(\ve{B}_t)=1,
  \label{eq:kdn2-gauge-fix}
\end{equation}
and compensating $\ve{P}_t$ so that $\ve{B}_t\otimes\ve{P}_t$ is unchanged.

At the covariance-family level, the nested special cases are
\begin{equation}
\begin{array}{ccl}
  \ve{B}_t=\bar b_t\ve{I}_V,\ \ve{R}_t=r_t\ve{I}_V
  &\Longrightarrow& \text{Diagonal KDN},\\
  \ve{P}_t=\bar p_t\ve{I}_K
  &\Longrightarrow& \text{value-wise isotropic KDN},\\
  \ve{B}_t=\bar b_t\ve{I}_V,\ \ve{P}_t=\bar p_t\ve{I}_K,\ \ve{R}_t=r_t\ve{I}_V
  &\Longrightarrow& \text{Isotropic KDN}.
\end{array}
  \label{eq:kdn2-special-cases}
\end{equation}
The scalar factor in each row can be absorbed into the other covariance factor before applying
the gauge in Equation~\eqref{eq:kdn2-gauge-fix}.  These are one-step covariance-family limits.
Recovering the complete one-axis recurrences additionally requires constraining the corresponding
factor and using its one-axis predictive projection; in particular, the channel-wise transition
in Equation~\eqref{eq:kdn2-factor-prediction} does not preserve
$\ve{P}_t\propto\ve{I}_K$ without the trace projection in
Equation~\eqref{eq:value-prediction}.

\subsection{Factor-preserving prediction}

The key-space mean transition remains
\begin{equation}
  \widehat{\ve{S}}_t=\ve{D}_t\ve{S}_{t-1},
  \qquad
  \ve{D}_t=\diag(\ve{\alpha}_t).
  \label{eq:kdn2-mean-prediction}
\end{equation}
A simple retained prediction updates the two covariance factors by
\begin{equation}
  \widehat{\ve{P}}_t
  =\ve{D}_t\ve{P}_{t-1}\ve{D}_t^{\tr}+\ve{\Omega}^{(k)}_t,
  \qquad
  \widehat{\ve{B}}_t
  =\ve{B}_{t-1}+\ve{\Omega}^{(v)}_t,
  \label{eq:kdn2-factor-prediction}
\end{equation}
with positive diagonal process factors.  This equation should be read as a factor-preserving
predictive parameterization.  It corresponds to the positive-semidefinite full process
covariance
\begin{equation}
\begin{aligned}
  \ve{Q}_t
  ={}&\ve{B}_{t-1}\otimes\ve{\Omega}^{(k)}_t
  +\ve{\Omega}^{(v)}_t\otimes
    (\ve{D}_t\ve{P}_{t-1}\ve{D}_t^{\tr})
  +\ve{\Omega}^{(v)}_t\otimes\ve{\Omega}^{(k)}_t,
\end{aligned}
  \label{eq:kdn2-implied-process-noise}
\end{equation}
because adding Equation~\eqref{eq:kdn2-implied-process-noise} to the transitioned prior gives
$\widehat{\ve{B}}_t\otimes\widehat{\ve{P}}_t$.  A generic additive separable process covariance
would instead produce a sum of Kronecker products and leave the retained family.

After prediction, apply the gauge normalization in Equation~\eqref{eq:kdn2-gauge-fix}.  This
normalization leaves the full predictive covariance and the full Kalman gain unchanged.  It does
rescale the factors used below: under
$(\widehat{\ve{B}}_t,\widehat{\ve{P}}_t)\mapsto
(\lambda\widehat{\ve{B}}_t,\widehat{\ve{P}}_t/\lambda)$,
$\ve{g}_t\mapsto\ve{g}_t/\lambda$ and
$\ve{\beta}_t\mapsto\lambda\ve{\beta}_t$, so their product is invariant.

\subsection{Exact innovation and conditional mean}

Define
\begin{equation}
  \ve{g}_t=\widehat{\ve{P}}_t\ve{k}_t,
  \qquad
  z_t=\ve{k}_t^{\tr}\widehat{\ve{P}}_t\ve{k}_t,
  \qquad
  \ve{\delta}_t=\ve{v}_t-\widehat{\ve{S}}_t^{\tr}\ve{k}_t.
  \label{eq:kdn2-key-statistics}
\end{equation}
The innovation covariance and state--observation cross-covariance are
\begin{equation}
  \ve{\Delta}_t
  =z_t\widehat{\ve{B}}_t+\ve{R}_t,
  \qquad
  \Cov(\ve{x}_t,\ve{v}_t\mid\mathcal{F}_{t-1})
  =\widehat{\ve{B}}_t\otimes\ve{g}_t.
  \label{eq:kdn2-joint-moments}
\end{equation}
The exact vectorized Kalman gain is therefore
\begin{equation}
  \ve{K}_t
  =(\widehat{\ve{B}}_t\otimes\ve{g}_t)\ve{\Delta}_t^{-1}.
  \label{eq:kdn2-vector-gain}
\end{equation}
When $\ve{R}_t=\diag(r_{t,1},\ldots,r_{t,V})$, define
\begin{equation}
  \beta_{t,j}
  =\frac{\widehat b_{t,j}}
    {r_{t,j}+\widehat b_{t,j}z_t},
  \qquad
  \ve{\beta}_t
  =(\beta_{t,1},\ldots,\beta_{t,V})^{\tr}.
  \label{eq:kdn2-value-gates}
\end{equation}
The exact posterior mean has the rank-one matrix update
\begin{equation}
  \ve{S}_t^{\star}
  =\widehat{\ve{S}}_t
   +\ve{g}_t(\ve{\beta}_t\odot\ve{\delta}_t)^{\tr}.
  \label{eq:kdn2-mean-update}
\end{equation}
Equivalently, grouping the predicted memory columnwise gives
\begin{equation}
\begin{aligned}
  \ve{s}_{t,j}^{\star}
  &=\left(
      \ve{I}_K-\beta_{t,j}\ve{g}_t\ve{k}_t^{\tr}
    \right)\widehat{\ve{s}}_{t,j}
    +\beta_{t,j}\ve{g}_tv_{t,j}
  \\
  &=\left(
      \ve{I}_K-\beta_{t,j}\ve{g}_t\ve{k}_t^{\tr}
    \right)\ve{D}_t\ve{s}_{t-1,j}
    +\beta_{t,j}\ve{g}_tv_{t,j},
  \qquad j=1,\ldots,V.
\end{aligned}
  \label{eq:kdn2-mean-column-factorized}
\end{equation}
This equation is the central behavioral generalization.  The key factor
$\widehat{\ve{P}}_t$ rotates and rescales the write direction through $\ve{g}_t$, while the value
factor $\widehat{\ve{B}}_t$ produces relative write strengths across value coordinates.  Their
absolute scales are gauge-dependent; only $\ve{g}_t\beta_{t,j}$ and ratios among the
$\beta_{t,j}$ are identifiable.  Setting
$\widehat{\ve{B}}_t=\ve{I}_V$ and $\ve{R}_t=r_t\ve{I}_V$ recovers the Diagonal-KDN gain;
setting $\widehat{\ve{P}}_t=\ve{I}_K$ recovers
Equation~\eqref{eq:value-memory-update}.

The one-step predictive negative log likelihood, up to the additive Gaussian constant, is
\begin{equation}
  \mathcal{L}_{\mathrm{pred},t}
  =\frac12\sum_{j=1}^{V}
    \left[
      \frac{\delta_{t,j}^2}{r_{t,j}+\widehat b_{t,j}z_t}
      +\log(r_{t,j}+\widehat b_{t,j}z_t)
    \right].
  \label{eq:kdn2-predictive-nll}
\end{equation}
Both the residual and its variance use the exclusive pre-write state.

\subsection{Exact posterior and failure of Kronecker closure}

The exact posterior covariance and precision are
\begin{equation}
\begin{aligned}
  \ve{\Sigma}_t^{\star}
  &={}
    \widehat{\ve{B}}_t\otimes\widehat{\ve{P}}_t
    -(\widehat{\ve{B}}_t\ve{\Delta}_t^{-1}\widehat{\ve{B}}_t)
      \otimes(\ve{g}_t\ve{g}_t^{\tr}),
  \\
  \ve{J}_t^{\star}
  &={}
    \widehat{\ve{B}}_t^{-1}\otimes\widehat{\ve{P}}_t^{-1}
    +\ve{R}_t^{-1}\otimes\ve{k}_t\ve{k}_t^{\tr}.
\end{aligned}
  \label{eq:kdn2-exact-posterior}
\end{equation}
For diagonal $\ve{R}_t$, column $j$ has covariance block
\begin{equation}
  \ve{\Sigma}_{t,j}^{\star}
  =\widehat b_{t,j}
   \left(
     \widehat{\ve{P}}_t
     -\beta_{t,j}\ve{g}_t\ve{g}_t^{\tr}
   \right).
  \label{eq:kdn2-posterior-block}
\end{equation}
Different $\beta_{t,j}$ values therefore produce differently shaped key-space covariance blocks,
not merely differently scaled copies of one common matrix.

\begin{proposition}[Exact separable closure]
Assume $K>1$, $\widehat{\ve{P}}_t\succ0$, and
$\ve{g}_t\ve{g}_t^{\tr}\ne\ve{0}$.  The blocks in
Equation~\eqref{eq:kdn2-posterior-block} can be written as
$b_{t,j}^{+}\ve{P}_t^{+}$ with one common, possibly dense, matrix $\ve{P}_t^{+}$ if and only if
all $\beta_{t,j}$ are equal.  For diagonal observation noise, this is equivalent to
\begin{equation}
  \frac{r_{t,j}}{\widehat b_{t,j}}=\rho_t
  \quad\text{for every }j,
  \qquad\text{or equivalently}\qquad
  \ve{R}_t=\rho_t\widehat{\ve{B}}_t.
  \label{eq:kdn2-compatible-noise}
\end{equation}
\end{proposition}

\begin{proof}
If $\beta_{t,j}=\beta_t$, then
$\ve{\Sigma}_{t,j}^{\star}=\widehat b_{t,j}
(\widehat{\ve{P}}_t-\beta_t\ve{g}_t\ve{g}_t^{\tr})$, which is separable.  Conversely, suppose
two blocks with coefficients $\beta_a\ne\beta_b$ are proportional.  Then for some $c>0$,
\[
  \widehat{\ve{P}}_t-\beta_a\ve{g}_t\ve{g}_t^{\tr}
  =c(\widehat{\ve{P}}_t-\beta_b\ve{g}_t\ve{g}_t^{\tr}).
\]
If $c\ne1$, the left rearrangement contains the full-rank matrix
$(1-c)\widehat{\ve{P}}_t$ while the other side has rank at most one, a contradiction.  Thus
$c=1$ and $\beta_a=\beta_b$.  Finally,
$\beta_{t,j}=1/(r_{t,j}/\widehat b_{t,j}+z_t)$, so equality of the gains is equivalent to
Equation~\eqref{eq:kdn2-compatible-noise}.
\end{proof}

The compatible choice $\ve{R}_t=\rho_t\widehat{\ve{B}}_t$ is mathematically closed in the value
factor, but it makes that factor cancel from the gain:
\begin{equation}
  \ve{S}_t^{\star}
  =\widehat{\ve{S}}_t
   +\frac{\ve{g}_t\ve{\delta}_t^{\tr}}{\rho_t+z_t}.
  \label{eq:kdn2-compatible-mean}
\end{equation}
The value uncertainty then has no effect on the posterior mean.  Moreover, the common key-space
posterior remains dense for a dense key, so the diagonal $\ve{P}_t$ restriction still requires
the usual KDN projection.  Useful value-dependent gains and exact Kronecker closure therefore
cannot generally be obtained simultaneously.

\subsection{Joint reverse-KL projection}

We retain the exact mean in Equation~\eqref{eq:kdn2-mean-update} and project only its covariance.
Let
\begin{equation}
  q_t
  =\Normal\!\left(
    \operatorname{vec}(\ve{S}_t^{\star}),
    \ve{B}_t\otimes\ve{P}_t
  \right),
  \qquad
  \ve{B}_t=\diag(\ve{b}_t),
  \quad
  \ve{P}_t=\diag(\ve{p}_t).
  \label{eq:kdn2-variational-family}
\end{equation}
Define four candidate-factor summaries
\begin{equation}
\begin{aligned}
  A_B&=\sum_{j=1}^{V}\frac{b_{t,j}}{\widehat b_{t,j}},
  &
  T_B&=\sum_{j=1}^{V}\frac{b_{t,j}}{r_{t,j}},
  \\
  A_P&=\sum_{i=1}^{K}\frac{p_{t,i}}{\widehat p_{t,i}},
  &
  Z_P&=\sum_{i=1}^{K}p_{t,i}k_{t,i}^2.
\end{aligned}
  \label{eq:kdn2-projection-summaries}
\end{equation}

\begin{proposition}[Joint reverse-KL projection]
Up to the factor gauge in Equation~\eqref{eq:kdn2-factor-gauge}, the stationary point of
$\KL(q_t\,\|\,\pi_t^{\star})$, where $\pi_t^{\star}$ has covariance
Equation~\eqref{eq:kdn2-exact-posterior}, satisfies
\begin{equation}
  b_{t,j}
  =\frac{K}{A_P/\widehat b_{t,j}+Z_P/r_{t,j}},
  \qquad
  p_{t,i}
  =\frac{V}{A_B/\widehat p_{t,i}+T_B k_{t,i}^2}.
  \label{eq:kdn2-reverse-kl-fixed-point}
\end{equation}
Equivalently, with
\begin{equation}
  L_{t,ij}
  =\frac{1}{\widehat p_{t,i}\widehat b_{t,j}}
   +\frac{k_{t,i}^2}{r_{t,j}},
  \qquad
  W_{t,ij}=p_{t,i}b_{t,j}L_{t,ij},
  \label{eq:kdn2-balancing-matrix}
\end{equation}
the positive matrix $\ve{W}_t$ has row sums $V$ and column sums $K$.
\end{proposition}

\begin{proof}
Using the precision in Equation~\eqref{eq:kdn2-exact-posterior}, the covariance-dependent part of
twice the reverse KL is
\begin{equation}
  A_BA_P+T_BZ_P-K\log\det\ve{B}_t-V\log\det\ve{P}_t
  +\mathrm{const}.
  \label{eq:kdn2-reverse-kl-objective}
\end{equation}
Differentiating with respect to $b_{t,j}$ and $p_{t,i}$ yields
Equation~\eqref{eq:kdn2-reverse-kl-fixed-point}.  Equivalently,
$\partial/\partial p_{t,i}=0$ gives
$\sum_jW_{t,ij}=V$, and $\partial/\partial b_{t,j}=0$ gives
$\sum_iW_{t,ij}=K$.  Multiplying every $b_{t,j}$ by $\lambda$ and dividing every $p_{t,i}$ by
$\lambda$ leaves both the objective and the full covariance unchanged, which is exactly the
factor gauge.
\end{proof}

In log-factor coordinates, Equation~\eqref{eq:kdn2-reverse-kl-objective} is convex and is strictly
convex after the gauge in Equation~\eqref{eq:kdn2-gauge-fix} is fixed.  The positive fixed point is
therefore unique in that gauge.  It is a matrix-scaling problem, and alternating updates
\begin{equation}
\begin{aligned}
  b_{t,j}&\leftarrow
  \frac{K}{A_P/\widehat b_{t,j}+Z_P/r_{t,j}},
  \\
  p_{t,i}&\leftarrow
  \frac{V}{A_B/\widehat p_{t,i}+T_Bk_{t,i}^2},
\end{aligned}
  \label{eq:kdn2-alternating-projection}
\end{equation}
followed by gauge normalization solve it without constructing a $K\times V$ matrix.  These are
the standard alternating matrix-scaling updates and converge to the unique gauge-equivalence
class because $L_{t,ij}>0$.  Each sweep costs $O(K+V)$ because $L_{t,ij}$ is the sum of two
rank-one terms, one strictly positive and one nonnegative.  The projection is nevertheless
coupled and iterative at every token.  It is not two independent M\"obius scans over time.

The one-sided boundaries recover the previous models.  Holding
$\ve{B}_t=\widehat{\ve{B}}_t$ fixed gives
\begin{equation}
  p_{t,i}
  =\left[
    \widehat p_{t,i}^{-1}
    +\left(\frac{1}{V}\sum_{j=1}^{V}\frac{\widehat b_{t,j}}{r_{t,j}}\right)k_{t,i}^2
  \right]^{-1}.
  \label{eq:kdn2-fixed-b-boundary}
\end{equation}
For $\widehat{\ve{B}}_t=\ve{I}_V$ and $\ve{R}_t=r_t\ve{I}_V$, this is the ordinary
Diagonal-KDN reverse-KL update.  Holding
$\ve{P}_t=\widehat{\ve{P}}_t$ fixed gives
\begin{equation}
  b_{t,j}
  =\left[
    \widehat b_{t,j}^{-1}
    +\frac{z_t}{K r_{t,j}}
  \right]^{-1},
  \label{eq:kdn2-fixed-p-boundary}
\end{equation}
which is the value-wise isotropic update when
$\widehat{\ve{P}}_t=\ve{I}_K$.  These equations describe one posterior projection.  Recovering
the full value-wise isotropic recurrence also requires the trace-projected prediction in
Equation~\eqref{eq:value-prediction}.  Applying both one-sided full-information updates
independently would double-count the same observation; the coupled projection is what allocates
the information between factors.

An optional protection scale may multiply the measurement term $k_{t,i}^2/r_{t,j}$ in
Equation~\eqref{eq:kdn2-balancing-matrix}.  If the same scale also tempers the likelihood and
changes the current gain, it defines a different Gaussian observation model.  If the exact gain
in Equation~\eqref{eq:kdn2-mean-update} is kept fixed, the scaled projection is instead a retained
future-protection calibration, as in the one-axis KDN variants; it is not the reverse-KL optimum
for Equation~\eqref{eq:observation-model}.

\subsection{A practical reference recurrence}

A sequential KDN-2 reference implementation can use the following token step:
\begin{enumerate}
  \item Predict $\widehat{\ve{S}}_t$, $\widehat{\ve{P}}_t$, and
    $\widehat{\ve{B}}_t$ with Equations~\eqref{eq:kdn2-mean-prediction} and
    \eqref{eq:kdn2-factor-prediction}.
  \item Fix the internal factor gauge by Equation~\eqref{eq:kdn2-gauge-fix}.
  \item Compute $\ve{g}_t$, $z_t$, $\ve{\beta}_t$, and the exact posterior mean
    $\ve{S}_t^{\star}$ from Equations~\eqref{eq:kdn2-key-statistics}--\eqref{eq:kdn2-mean-update}.
  \item Initialize the covariance projection at the predictive factors and run a fixed number of
    alternating updates in Equation~\eqref{eq:kdn2-alternating-projection}.
  \item Reapply the gauge normalization and carry the projected factors to token $t+1$.
\end{enumerate}
With $N_{\mathrm{proj}}$ alternating sweeps, uncertainty work is
$O(N_{\mathrm{proj}}(K+V))$ per token and the memory update is $O(KV)$.  This is suitable as a
correctness oracle and as a small-scale experiment.  It is not yet a drop-in replacement for the
current scan-parallel KDN kernel.

\section{Connection to Gated DeltaNet-2}
\label{sec:gdn2-connection}

\subsection{The two updates side by side}

Use the same predicted memory
$\widehat{\ve{S}}_t=\ve{D}_t\ve{S}_{t-1}$.  Gated DeltaNet-2 (GDN-2) can be written as
\begin{equation}
\begin{aligned}
  \ve{S}_t^{\mathrm{GDN2}}
  &={}
    \widehat{\ve{S}}_t
    -\ve{k}_t
      (\ve{b}_t^{\mathrm{erase}}\odot\ve{k}_t)^{\tr}\widehat{\ve{S}}_t
    +\ve{k}_t(\ve{w}_t\odot\ve{v}_t)^{\tr},
  \\
  \ve{b}_t^{\mathrm{erase}}&\in\R^K,
  \qquad
  \ve{w}_t\in\R^V,
\end{aligned}
  \label{eq:gdn2-update}
\end{equation}
where $\ve{b}_t^{\mathrm{erase}}$ is a key-channel erase gate and $\ve{w}_t$ is a value-channel
write gate.  Both are predicted directly from the current token and are sigmoid-valued in the
default implementation; an optional GDN-2 variant expands only the erase-gate range.  The keys are
$\ell_2$-normalized in both model families.  Defining the pseudo-key and pseudo-value
\begin{equation}
  \widetilde{\ve{k}}_t=\ve{b}_t^{\mathrm{erase}}\odot\ve{k}_t,
  \qquad
  \widetilde{\ve{v}}_t=\ve{w}_t\odot\ve{v}_t,
  \label{eq:gdn2-pseudo-observation}
\end{equation}
gives the residual-like form
\begin{equation}
  \ve{S}_t^{\mathrm{GDN2}}
  =\widehat{\ve{S}}_t
   +\ve{k}_t
    \left(
      \widetilde{\ve{v}}_t
      -\widehat{\ve{S}}_t^{\tr}\widetilde{\ve{k}}_t
    \right)^{\tr}.
  \label{eq:gdn2-pseudo-residual}
\end{equation}
This resembles a Kalman residual update algebraically, but its erased read uses
$\widetilde{\ve{k}}_t$ while its injected target is $\widetilde{\ve{v}}_t$.  In general it is not
the posterior mean for the original observation
$\ve{v}_t=\widetilde{\ve{S}}_t^{\tr}\ve{k}_t+\ve{e}^{\mathrm{obs}}_t$.

Expanding the KDN-2 mean update in Equation~\eqref{eq:kdn2-mean-update} gives
\begin{equation}
  \ve{S}_t^{\mathrm{KDN2}}
  =\widehat{\ve{S}}_t
   -\ve{g}_t\ve{k}_t^{\tr}\widehat{\ve{S}}_t\diag(\ve{\beta}_t)
   +\ve{g}_t(\ve{\beta}_t\odot\ve{v}_t)^{\tr},
  \qquad
  \ve{g}_t=\widehat{\ve{P}}_t\ve{k}_t.
  \label{eq:kdn2-expanded-gdn2-comparison}
\end{equation}
Both models therefore use an $O(KV)$ low-rank correction, and both can modulate the incoming value
coordinatewise.  Their factorizations differ:
\begin{center}
\begin{tabular}{@{}llll@{}}
  \toprule
  Model & outer write direction & erased read direction & value multiplier \\
  \midrule
  GDN-2 & $\ve{k}_t$ & $\ve{b}_t^{\mathrm{erase}}\odot\ve{k}_t$ & $\ve{w}_t$ \\
  KDN-2 & $\ve{g}_t=\widehat{\ve{P}}_t\ve{k}_t$ & $\ve{k}_t$ & $\ve{\beta}_t$ \\
  \bottomrule
\end{tabular}
\end{center}
For KDN-2, the same $\beta_{t,j}$ multiplies the observation write and the erasure of the predicted
value in column $j$.  This tie is forced by the Kalman innovation.  GDN-2 instead predicts
$\ve{b}_t^{\mathrm{erase}}$ and $\ve{w}_t$ independently, deliberately decoupling how much old
content is erased from how much new value is injected.

The two covariance factors explain the two visible parts of the KDN-2 update.  The key factor
$\widehat{\ve{P}}_t$ changes the outer write direction from $\ve{k}_t$ to $\ve{g}_t$, while the
value factor $\widehat{\ve{B}}_t$ and observation noise determine the relative entries of
$\ve{\beta}_t$.  These gates are functions of the carried uncertainty state and therefore depend
on the history, not only on the current token representation.

\subsection{Exact intersection of the update families}

The similarity between Equations~\eqref{eq:gdn2-update} and
\eqref{eq:kdn2-expanded-gdn2-comparison} does not imply that either family contains the other.
Their common intersection is much smaller.

\begin{proposition}[Algebraic GDN-2--KDN-2 intersection]
At the level of the memory-update equations, assume a nondegenerate update with
$\ve{k}_t\ne\ve{0}$ and require
Equations~\eqref{eq:gdn2-update} and \eqref{eq:kdn2-expanded-gdn2-comparison} to agree for arbitrary
$\widehat{\ve{S}}_t$ and $\ve{v}_t$.  Then there must be scalars $\lambda_t>0$ and
$\beta_t>0$ such that
\begin{equation}
\begin{aligned}
  \ve{g}_t&=\lambda_t\ve{k}_t,
  &
  \ve{\beta}_t&=\beta_t\ve{1}_V,
  \\
  \ve{w}_t&=\lambda_t\beta_t\ve{1}_V,
  &
  \ve{b}_t^{\mathrm{erase}}\odot\ve{k}_t
  &=\lambda_t\beta_t\ve{k}_t.
\end{aligned}
  \label{eq:gdn2-kdn2-intersection}
\end{equation}
Conversely, these conditions make the two updates identical.
\end{proposition}

\begin{proof}
Equality of the coefficients multiplying each arbitrary $v_{t,j}$ requires
\begin{equation}
  \beta_{t,j}\ve{g}_t=w_{t,j}\ve{k}_t.
  \label{eq:gdn2-kdn2-write-condition}
\end{equation}
For a nonzero write this forces $\ve{g}_t=\lambda_t\ve{k}_t$ and
$w_{t,j}=\lambda_t\beta_{t,j}$.  Equality of the linear maps applied to each arbitrary memory
column then requires
\begin{equation}
  \lambda_t\beta_{t,j}\ve{k}_t\ve{k}_t^{\tr}
  =\ve{k}_t(\ve{b}_t^{\mathrm{erase}}\odot\ve{k}_t)^{\tr}
  \qquad\text{for every }j.
  \label{eq:gdn2-kdn2-erase-condition}
\end{equation}
The right-hand side is shared across value columns, so every $\beta_{t,j}$ must equal one scalar
$\beta_t$, and
$\ve{b}_t^{\mathrm{erase}}\odot\ve{k}_t
=\lambda_t\beta_t\ve{k}_t$.  Substitution proves sufficiency.
\end{proof}

For the implemented default GDN-2 parameterization, the converse additionally requires
$\lambda_t\beta_t\in(0,1)$ so that both gates lie in their sigmoid ranges.  The optional
negative-eigenvalue variant permits an erase coefficient up to two but leaves the write gate in
$(0,1)$.

For a dense key and diagonal $\widehat{\ve{P}}_t$, the condition
$\ve{g}_t=\lambda_t\ve{k}_t$ requires
$\widehat p_{t,i}=\lambda_t$ on the support of $\ve{k}_t$.  The intersection therefore reduces to
the scalar-gated delta update
\begin{equation}
  \ve{S}_t
  =\widehat{\ve{S}}_t
   +\gamma_t\ve{k}_t
    (\ve{v}_t-\widehat{\ve{S}}_t^{\tr}\ve{k}_t)^{\tr},
  \qquad
  \gamma_t=\lambda_t\beta_t,
  \label{eq:gdn2-kdn2-scalar-intersection}
\end{equation}
which is the KDA form.  Away from this intersection, KDN-2 can rotate the write direction through
key-space uncertainty, which GDN-2 cannot express with its fixed outer direction $\ve{k}_t$;
GDN-2 can independently choose erase and write gates, which a coherent KDN-2 innovation cannot
express.

The one-axis boundaries expose the two separate mismatches.  When
$\widehat{\ve{P}}_t\propto\ve{I}_K$ but $\ve{\beta}_t$ remains value dependent, KDN-2 reduces to
the value-wise isotropic update: it shares GDN-2's outer direction and value-wise write scaling,
but still uses the same $\beta_{t,j}$ for both halves of each column's innovation.  When
$\widehat{\ve{B}}_t\propto\ve{I}_V$ and $\ve{R}_t\propto\ve{I}_V$, KDN-2 reduces to Diagonal KDN:
its value gate becomes common, but the anisotropic direction
$\widehat{\ve{P}}_t\ve{k}_t$ generally cannot be represented by GDN-2.

\subsection{Interpretation and controlled comparisons}

GDN-2 may be viewed as a discriminative relaxation of an uncertainty-gated delta rule.  It
predicts the two sides of the update directly and gains token-local flexibility, but it does not
carry a covariance whose predictive variance or posterior contraction constrains those gates.
KDN-2 instead derives its update from the separable predictive covariance and diagonal observation
noise.  Its restrictions provide a calibrated predictive distribution and a causal notion of
future protection, at the cost of the coupled covariance projection derived above.

The connection suggests four controlled ablations:
\begin{enumerate}
  \item \textbf{Native scalar innovation tying.}  Within the existing GDN-2 operator, constrain
    $\ve{b}_t^{\mathrm{erase}}=\gamma_t\ve{1}_K$ and
    $\ve{w}_t=\gamma_t\ve{1}_V$.  This tests the benefit of decoupling while retaining the native
    shared erase system.
  \item \textbf{Value-wise tied innovation.}  Replace the shared erase operator by
    $-\ve{k}_t\ve{k}_t^{\tr}\widehat{\ve{S}}_t\diag(\ve{\beta}_t)$ and use the same
    $\ve{\beta}_t$ on the incoming value.  This is a structural recurrence ablation, not merely a
    gate constraint on GDN-2.
  \item \textbf{Gain direction.}  Replace $\ve{g}_t=\widehat{\ve{P}}_t\ve{k}_t$ by a collinear
    direction
    \begin{equation*}
      \ve{g}_{t,\parallel}
      =\frac{\ve{k}_t^{\tr}\ve{g}_t}{\|\ve{k}_t\|_2^2}\ve{k}_t,
    \end{equation*}
    which preserves $\ve{k}_t^{\tr}\ve{g}_t$ and hence the current-key contraction while
    removing the off-key component.
  \item \textbf{Stateful versus direct gates.}  Match instantaneous write strengths while
    comparing uncertainty-derived gates with gates predicted only from the current token.
\end{enumerate}
Useful diagnostics are the cosine between $\ve{g}_t$ and $\ve{k}_t$, dispersion across
$\beta_{t,j}$, disagreement between GDN-2 erase and write strengths, predictive innovation NLL,
and repeated-key protection.

\section{Parallelism and approximation choices}
\label{sec:parallelism}

The value-wise model in Section~\ref{sec:value-isotropic} has independent scalar information
recurrences, but its value-dependent gains require separate memory systems.  KDN-2 adds a second
difficulty: its covariance projection couples all key and value factors at every token.  Three
implementation levels should be distinguished:
\begin{enumerate}
  \item \textbf{Sequential joint projection.}  Run the exact gain and several alternating
    reverse-KL sweeps per token.  This is the reference semantics.
  \item \textbf{Grouped values.}  Share one $b$ and one gain across a small value group.  This
    reduces the number of memory systems while retaining coarse value reliability.
  \item \textbf{Chunk-frozen factors.}  Freeze or lag one factor within a chunk, use the
    corresponding one-sided update, and commit a joint projection at the boundary.  This may
    recover useful matrix operations, but it is an approximation whose error must be measured
    against the sequential reference.
\end{enumerate}
The compatible-noise construction in Equation~\eqref{eq:kdn2-compatible-noise} restores a shared
gain, but it also removes the value factor from the mean and is therefore a control rather than a
full KDN-2 implementation.
Appendix~\ref{app:value-wise-chunkwise} derives the exact value-batched compact-WY system for the
value-wise model and makes explicit where its computation ceases to be shared across value
coordinates.
Appendix~\ref{app:kdn2-chunkwise} then derives an exact two-stage implementation of two-axis
KDN-2: a causal coupled-factor prepass followed by a value-batched chunkwise mean pass.

\section{Parameterization, gauges, and validation}
\label{sec:validation}

\paragraph{Positive factors.}
Let $\ve{h}_t$ denote the observed token representation supplied to the uncertainty heads.
Predict diagonal process increments and observation noise through positive floors, for example
\begin{equation}
  \ve{\omega}^{(k)}_t
  =\omega_{k,\min}+\operatorname{softplus}(f_k(\ve{h}_t)),
  \quad
  \ve{\omega}^{(v)}_t
  =\omega_{v,\min}+\operatorname{softplus}(f_v(\ve{h}_t)),
  \quad
  \ve{r}_t=r_{\min}+\operatorname{softplus}(f_r(\ve{h}_t)).
  \label{eq:kdn2-positive-parameterization}
\end{equation}
Fixing $\ve{r}_t=r\ve{1}_V$ while learning value-specific
$\ve{\omega}^{(v)}_t$ is a useful first experiment: it breaks value symmetry without learning the
most direct $b_{t,j}/r_{t,j}$ confound.  In addition to the exact internal factor gauge in
Equation~\eqref{eq:kdn2-factor-gauge}, jointly scaling the full covariance and observation noise
produces the usual Kalman covariance-scale gauge.

\paragraph{Symmetry breaking.}
Initialize $\ve{B}_0=\ve{I}_V$ and $\ve{P}_0=\ve{I}_K$ for a neutral gauge.  Value uncertainty
remains tied unless the initial value factor, value-process uncertainty, or observation noise
varies across value coordinates.  Realized values alone cannot change covariance symmetry.

\paragraph{Required tests.}
A KDN-2 implementation should pass the following checks before language-model training:
\begin{enumerate}
  \item Compare the exact mean and covariance in
    Equations~\eqref{eq:kdn2-mean-update} and \eqref{eq:kdn2-exact-posterior} against a dense
    Kalman update on small random systems.
  \item Verify the Diagonal-KDN and value-wise isotropic limits in
    Equations~\eqref{eq:kdn2-fixed-b-boundary} and \eqref{eq:kdn2-fixed-p-boundary}.
  \item Check factor-gauge invariance numerically under
    $(\ve{B},\ve{P})\mapsto(\lambda\ve{B},\ve{P}/\lambda)$.
  \item Confirm that the alternating projection satisfies the row- and column-sum conditions in
    Equation~\eqref{eq:kdn2-balancing-matrix} and decreases the reverse-KL objective.
  \item Compare zero, one, two, and converged projection sweeps; report covariance error, gain
    error at later tokens, throughput, and peak memory.
  \item Test repeated keys, orthogonal keys with identical elementwise squares, value-specific
    corruption, and value-basis rotations.  The last test exposes the basis dependence of a
    diagonal value factor.
\end{enumerate}

\section{Limitations and experimental contract}
\label{sec:limitations}

KDN-2 represents a rank-one factorization of the $K\times V$ table of entrywise variances.  It is
strictly richer than either one-axis family but much smaller than a fully independent diagonal
covariance with $KV$ states.  It still assumes no cross-key or cross-value posterior covariance,
depends on the learned key and value bases, and introduces both an exact factor gauge and strong
confounding between value variance and observation noise.

The most important limitation is computational rather than notational.  The exact posterior mean
is inexpensive, but the joint covariance projection is coupled and the value-dependent gains
break the shared memory transition.  A claimed KDN-2 kernel must therefore specify which
projection it implements, how many alternating sweeps it uses, how the gauge is fixed, and whether
the memory system is value-wise, grouped, or chunk-approximated.  Results should be compared with
the two one-sided boundaries at matched parameter count and effective write strength.

The decisive initial experiment is a small synthetic ablation with four arms:
\begin{equation}
\begin{aligned}
  &\text{Diagonal KDN},
  &&\text{value-wise isotropic KDN},\\
  &\text{KDN-2 with one projection sweep},
  &&\text{KDN-2 with a converged projection}.
\end{aligned}
  \label{eq:kdn2-initial-ablation}
\end{equation}
This separates the benefit of two-axis uncertainty from the cost and approximation error of the
joint projection before committing to a production kernel.

\section{Conclusion}

Value-wise isotropic uncertainty and Diagonal KDN allocate the same kind of scalar uncertainty
budget to opposite axes of the memory matrix.  KDN-2 combines them through the separable
covariance $\ve{B}_t\otimes\ve{P}_t$.  Its exact mean update cleanly factors into a key-wise
uncertainty direction and value-wise uncertainty gates.  Its covariance update does not factor
cleanly: useful value-dependent gains destroy Kronecker closure, and the principled reverse-KL
projection is a coupled matrix-balancing problem.  This distinction is the central design
constraint.  KDN-2 is mathematically well defined as a sequential filter and a useful research
model, but a scan-parallel approximation must be stated and validated separately.

\appendix

\section{Chunkwise implementation of value-wise isotropic KDN}
\label{app:value-wise-chunkwise}

This appendix derives an exact chunkwise evaluation of the value-wise isotropic model in
Section~\ref{sec:value-isotropic}.  The derivation follows the compact-WY construction for the
delta rule: first remove the diagonal decay, then expose the causal residuals as a unit-lower
triangular system, and finally express all reads and the outgoing state as matrix products.  The
only structural change is important: every value coordinate has its own gain sequence, so the
triangular system is value batched rather than shared.

\subsection{Two-stage scan}
\label{app:value-wise-two-stage}

The uncertainty state does not depend on the realized values or on the retained memory.  Hence it
can be evaluated before the memory recurrence.  For each value coordinate $j$, write the
information state homogeneously as
\begin{equation}
  \ve{z}^{(c)}_{t,j}
  =\begin{bmatrix}n_{t,j}\\d_{t,j}\end{bmatrix},
  \qquad
  c_{t,j}=\frac{n_{t,j}}{d_{t,j}},
  \qquad
  \ve{z}^{(c)}_{t,j}=\ve{M}_{t,j}\ve{z}^{(c)}_{t-1,j},
  \label{eq:value-chunk-homogeneous-information}
\end{equation}
where $\ve{M}_{t,j}$ is the $2\times2$ matrix in
Equation~\eqref{eq:value-mobius}.  Products of these matrices are associative.  A chunked
reduce--scan--replay pass therefore produces every exclusive $c_{t-1,j}$ and hence every
$\widehat b_{t,j}$ and $\beta_{t,j}$ before the memory pass.  Explicitly, the exclusive state
starts from $\ve{z}^{(c)}_{0,j}=[c_{0,j},1]^{\tr}$ and emits
\begin{equation}
  \widehat c_{t,j}
  =\frac{c_{t-1,j}}{a_t+\omega_{t,j}c_{t-1,j}},
  \qquad
  \beta_{t,j}
  =\frac{1}{r_{t,j}\widehat c_{t,j}+\|\ve{k}_t\|_2^2},
  \label{eq:value-chunk-exclusive-gain}
\end{equation}
before $\ve{M}_{t,j}$ incorporates token $t$.  In homogeneous coordinates, the same emission can
avoid intermediate divisions: with
$\chi_{t,j}=a_td_{t-1,j}+\omega_{t,j}n_{t-1,j}$,
\begin{equation}
  \widehat c_{t,j}=\frac{n_{t-1,j}}{\chi_{t,j}},
  \qquad
  \beta_{t,j}
  =\frac{\chi_{t,j}}
   {r_{t,j}n_{t-1,j}+\|\ve{k}_t\|_2^2\chi_{t,j}}.
  \label{eq:value-chunk-homogeneous-gain}
\end{equation}
Homogeneous vectors and chunk summaries may be rescaled after multiplication without changing
their represented precision.  The gain and memory chunk lengths need not be the same.  What
follows conditions on these precomputed gains.

Consider one memory chunk of length $C$, use local indices $r=1,\ldots,C$, and suppress its global
chunk index.  Let $\ve{S}^{\mathrm{in}}=\ve{S}_0$ be the incoming memory.  Substituting
$\widehat{\ve{S}}_r=\ve{D}_r\ve{S}_{r-1}$ into
Equation~\eqref{eq:value-memory-update} gives
\begin{equation}
  \ve{S}_r
  =\ve{D}_r\ve{S}_{r-1}
   +\ve{k}_r
    \left[
      \ve{\beta}_r\odot
      \left(
        \ve{v}_r-(\ve{D}_r\ve{S}_{r-1})^{\tr}\ve{k}_r
      \right)
    \right]^{\tr}.
  \label{eq:value-chunk-original-recurrence}
\end{equation}

\subsection{Removing the diagonal decay}
\label{app:value-wise-decay-normalization}

Assume the gated transition has $\alpha_{r,i}>0$.  Define the within-chunk cumulative decay
\begin{equation}
  \ve{\gamma}_r=\prod_{\ell=1}^{r}\ve{\alpha}_\ell,
  \qquad
  \ve{\Gamma}_r=\diag(\ve{\gamma}_r),
  \qquad
  \ve{\gamma}_0=\ve{1},
  \label{eq:value-chunk-cumulative-decay}
\end{equation}
where the product is elementwise.  Since
$\ve{D}_r\ve{\Gamma}_{r-1}=\ve{\Gamma}_r$, write
\begin{equation}
  \ve{S}_r=\ve{\Gamma}_r\ve{Y}_r,
  \qquad
  \overline{\ve{k}}_r=\ve{\Gamma}_r\ve{k}_r
  =\ve{\gamma}_r\odot\ve{k}_r,
  \qquad
  \overline{\ve{p}}_r=\ve{\Gamma}_r^{-1}\ve{k}_r
  =\ve{k}_r\oslash\ve{\gamma}_r.
  \label{eq:value-chunk-normalized-keys}
\end{equation}
The incoming normalized state is $\ve{Y}_0=\ve{S}^{\mathrm{in}}$.  Multiplying
Equation~\eqref{eq:value-chunk-original-recurrence} by $\ve{\Gamma}_r^{-1}$ gives the undecayed
two-key delta rule
\begin{equation}
  \ve{Y}_r
  =\ve{Y}_{r-1}
   +\overline{\ve{p}}_r
    \left[
      \ve{\beta}_r\odot
      \left(
        \ve{v}_r-\ve{Y}_{r-1}^{\tr}\overline{\ve{k}}_r
      \right)
    \right]^{\tr}.
  \label{eq:value-chunk-normalized-recurrence}
\end{equation}
Thus $\overline{\ve{k}}_r$ is the read key and
$\beta_{r,j}\overline{\ve{p}}_r$ is the write key for value coordinate $j$.  Implementations
should form the required decay ratios from cumulative log decays rather than explicitly dividing
by a potentially small $\ve{\gamma}_r$.

\subsection{Causal innovations as value-batched triangular systems}
\label{app:value-wise-triangular-system}

Stack the normalized read and base-write keys by row,
\begin{equation}
  [\overline{\ve{K}}]_{r,:}=\overline{\ve{k}}_r^{\tr},
  \qquad
  [\overline{\ve{P}}]_{r,:}=\overline{\ve{p}}_r^{\tr},
  \qquad
  \ve{L}=\operatorname{tril}
  \left(\overline{\ve{K}}\overline{\ve{P}}^{\tr},-1\right).
  \label{eq:value-chunk-shared-score}
\end{equation}
The strictly lower matrix $\ve{L}\in\R^{C\times C}$ is shared by all value coordinates.  For
coordinate $j$, define
\begin{equation}
\begin{aligned}
  \ve{v}_j&=(v_{1,j},\ldots,v_{C,j})^{\tr},
  &
  \ve{\Lambda}^{(\beta)}_j
  &=\diag(\beta_{1,j},\ldots,\beta_{C,j}),
  \\
  \varepsilon_{r,j}
  &=v_{r,j}-\overline{\ve{k}}_r^{\tr}\ve{y}_{r-1,j},
  &
  \ve{\varepsilon}_j
  &=(\varepsilon_{1,j},\ldots,\varepsilon_{C,j})^{\tr},
  \\
  \ve{\xi}_j&=\ve{\Lambda}^{(\beta)}_j\ve{\varepsilon}_j.
\end{aligned}
  \label{eq:value-chunk-innovation-definitions}
\end{equation}
Unrolling only the preceding writes in
Equation~\eqref{eq:value-chunk-normalized-recurrence} yields
\begin{equation}
  \ve{\varepsilon}_j
  =\ve{v}_j-\overline{\ve{K}}\ve{s}^{\mathrm{in}}_j-\ve{L}\ve{\xi}_j,
  \qquad
  \ve{\xi}_j=\ve{\Lambda}^{(\beta)}_j\ve{\varepsilon}_j.
  \label{eq:value-chunk-causal-unrolling}
\end{equation}
Consequently, the complete set of causal innovations in the chunk is the unique solution of either
equivalent unit-lower system
\begin{equation}
\begin{aligned}
  (\ve{I}_C+\ve{L}\ve{\Lambda}^{(\beta)}_j)\ve{\varepsilon}_j
  &=\ve{v}_j-\overline{\ve{K}}\ve{s}^{\mathrm{in}}_j,
  \\
  (\ve{I}_C+\ve{\Lambda}^{(\beta)}_j\ve{L})\ve{\xi}_j
  &=\ve{\Lambda}^{(\beta)}_j
    (\ve{v}_j-\overline{\ve{K}}\ve{s}^{\mathrm{in}}_j).
\end{aligned}
  \label{eq:value-chunk-two-triangular-forms}
\end{equation}
The first form remains well defined when a gain is zero and is preferable to writing a system
with $(\ve{\Lambda}^{(\beta)}_j)^{-1}$.
Let $\ve{\Xi}=[\ve{\xi}_1,\ldots,\ve{\xi}_V]\in\R^{C\times V}$, with
$\ve{\xi}_r^{\tr}$ denoting row $r$.  Equivalently, all value coordinates can be advanced together
by the fused forward substitution
\begin{equation}
  \ve{\xi}_r
  =\ve{\beta}_r\odot
   \left(
     \ve{v}_r-(\ve{S}^{\mathrm{in}})^{\tr}\overline{\ve{k}}_r
     -\sum_{i<r}L_{ri}\ve{\xi}_i
   \right),
  \qquad
  \ve{\xi}_r\in\R^V.
  \label{eq:value-chunk-fused-substitution}
\end{equation}
This form makes the reusable geometry $\ve{L}$ and the elementwise value-axis gain explicit.

To recover the usual compact-WY notation, define
\begin{equation}
  \ve{A}_j=(\ve{I}_C+\ve{L}\ve{\Lambda}^{(\beta)}_j)^{-1},
  \qquad
  \ve{\eta}_j=\ve{A}_j\ve{v}_j,
  \qquad
  \ve{W}_j=\ve{A}_j\overline{\ve{K}}.
  \label{eq:value-chunk-wy-factors}
\end{equation}
Then
\begin{equation}
  \ve{\varepsilon}_j
  =\ve{\eta}_j-\ve{W}_j\ve{s}^{\mathrm{in}}_j,
  \qquad
  \ve{\xi}_j
  =\ve{\Lambda}^{(\beta)}_j
   (\ve{\eta}_j-\ve{W}_j\ve{s}^{\mathrm{in}}_j).
  \label{eq:value-chunk-wy-residual}
\end{equation}
The inverse in Equation~\eqref{eq:value-chunk-wy-factors} is notation: an implementation uses a
unit-lower-triangular solve.  The factors $\ve{\eta}_j$ and $\ve{W}_j$ depend only on inputs from the
current chunk and may therefore be prepared for all chunks in parallel, before the incoming chunk
states are known.

The connection to the compact-WY product is explicit.  Define the normalized rank-one transition
\begin{equation}
  \ve{E}_{r,j}
  =\ve{I}_K
   -\beta_{r,j}\overline{\ve{p}}_r\overline{\ve{k}}_r^{\tr}.
  \label{eq:value-chunk-rank-one-transition}
\end{equation}
Then the product of all within-chunk transitions is
\begin{equation}
  \ve{E}_{C,j}\ve{E}_{C-1,j}\cdots\ve{E}_{1,j}
  =\ve{I}_K
   -\overline{\ve{P}}^{\tr}
    \ve{\Lambda}^{(\beta)}_j\ve{W}_j.
  \label{eq:value-chunk-compact-wy-identity}
\end{equation}
The same transformed residuals also collect all forced writes, so
\begin{equation}
  \ve{y}_{C,j}
  =\left(
     \ve{I}_K
     -\overline{\ve{P}}^{\tr}
      \ve{\Lambda}^{(\beta)}_j\ve{W}_j
   \right)\ve{s}^{\mathrm{in}}_j
   +\overline{\ve{P}}^{\tr}
    \ve{\Lambda}^{(\beta)}_j\ve{\eta}_j
  =\ve{s}^{\mathrm{in}}_j+\overline{\ve{P}}^{\tr}\ve{\xi}_j.
  \label{eq:value-chunk-compact-wy-state}
\end{equation}

\subsection{Parallel reads and the outgoing state}
\label{app:value-wise-chunk-output}

The normalized state after any local prefix $r$ is
\begin{equation}
  \ve{Y}_r
  =\ve{S}^{\mathrm{in}}
   +\sum_{i=1}^{r}\overline{\ve{p}}_i\ve{\xi}_i^{\tr},
  \label{eq:value-chunk-prefix-state}
\end{equation}
where $\ve{\xi}_i^{\tr}$ is row $i$ of $\ve{\Xi}$.  In particular, the outgoing physical memory is
\begin{equation}
  \ve{S}^{\mathrm{out}}
  =\ve{\Gamma}_C
   \left(
     \ve{S}^{\mathrm{in}}+\overline{\ve{P}}^{\tr}\ve{\Xi}
   \right).
  \label{eq:value-chunk-outgoing-state}
\end{equation}

For completeness, suppose the mixer reads after the current write with queries
$\ve{q}_r\in\R^K$ and outputs $\ve{o}_r=\ve{S}_r^{\tr}\ve{q}_r$.  Define
\begin{equation}
  [\overline{\ve{Q}}]_{r,:}
  =(\ve{\Gamma}_r\ve{q}_r)^{\tr},
  \qquad
  \ve{L}_{qp}
  =\operatorname{tril}
   (\overline{\ve{Q}}\overline{\ve{P}}^{\tr}).
  \label{eq:value-chunk-query-scores}
\end{equation}
The diagonal is retained because the read is post-write.  All outputs in the chunk are
\begin{equation}
  \ve{O}
  =\overline{\ve{Q}}\ve{S}^{\mathrm{in}}
   +\ve{L}_{qp}\ve{\Xi}.
  \label{eq:value-chunk-outputs}
\end{equation}
A pre-write read instead uses the strictly lower part of the same score matrix.

The inverse cumulative decay in $\overline{\ve{P}}$ need not be materialized.  The entries used in
the shared score matrices are
\begin{equation}
\begin{aligned}
  L_{ri}
  &=\ve{k}_r^{\tr}
    \diag(\ve{\gamma}_r\oslash\ve{\gamma}_i)\ve{k}_i,
    \qquad i<r,
  \\
  (L_{qp})_{ri}
  &=\ve{q}_r^{\tr}
    \diag(\ve{\gamma}_r\oslash\ve{\gamma}_i)\ve{k}_i,
    \qquad i\le r.
\end{aligned}
  \label{eq:value-chunk-decay-ratio-scores}
\end{equation}
Likewise, define the tail-adjusted write matrix by
\begin{equation}
  [\ve{K}_{\mathrm{tail}}]_{r,:}
  =\left[
    (\ve{\gamma}_C\oslash\ve{\gamma}_r)\odot\ve{k}_r
   \right]^{\tr}.
  \label{eq:value-chunk-tail-key}
\end{equation}
Equation~\eqref{eq:value-chunk-outgoing-state} can then be evaluated as
\begin{equation}
  \ve{S}^{\mathrm{out}}
  =\ve{\Gamma}_C\ve{S}^{\mathrm{in}}
   +\ve{K}_{\mathrm{tail}}^{\tr}\ve{\Xi}.
  \label{eq:value-chunk-stable-outgoing-state}
\end{equation}
The ratios in Equations~\eqref{eq:value-chunk-decay-ratio-scores} and
\eqref{eq:value-chunk-tail-key} contain only the decay between the write and the read or chunk
boundary.  They should be computed from cumulative log-decay differences.  This remains stable
even when the individual cumulative products underflow.

\begin{proposition}[Exactness of the value-batched chunk form]
For fixed gains $\{\ve{\beta}_r\}_{r=1}^{C}$ and positive diagonal decays, the solutions of
Equation~\eqref{eq:value-chunk-two-triangular-forms} and the prefix reconstruction in
Equation~\eqref{eq:value-chunk-prefix-state} equal the token-by-token recurrence in
Equation~\eqref{eq:value-chunk-original-recurrence}.  Consequently,
Equations~\eqref{eq:value-chunk-outgoing-state} and \eqref{eq:value-chunk-outputs} give its exact
chunk boundary and post-write outputs.
\end{proposition}

\begin{proof}
Equation~\eqref{eq:value-chunk-normalized-recurrence} implies
$\ve{y}_{r-1,j}=\ve{s}^{\mathrm{in}}_j+
\sum_{i<r}\overline{\ve{p}}_i\xi_{i,j}$.  Substitution into the definition of
$\varepsilon_{r,j}$ gives row $r$ of Equation~\eqref{eq:value-chunk-causal-unrolling}.  Its
coefficient matrix is unit lower triangular, so its solution is unique and equals the sequential
innovations.  Substituting those innovations back into the normalized recurrence gives
Equation~\eqref{eq:value-chunk-prefix-state}; rescaling and applying the queries gives
Equations~\eqref{eq:value-chunk-outgoing-state} and \eqref{eq:value-chunk-outputs}.
\end{proof}

\subsection{Parallel schedule, cost, and sharing limit}
\label{app:value-wise-chunk-cost}

An exact implementation has the following schedule:
\begin{enumerate}
  \item Compute all $\beta_{t,j}$ with the independent M\"obius scans in
    Section~\ref{app:value-wise-two-stage}.
  \item For every memory chunk in parallel, form the cumulative decay, the normalized keys, and
    the shared score matrices $\ve{L}$ and $\ve{L}_{qp}$.
  \item Solve Equation~\eqref{eq:value-chunk-two-triangular-forms} batched over value coordinates,
    then evaluate Equations~\eqref{eq:value-chunk-outputs} and
    \eqref{eq:value-chunk-outgoing-state}.  Only the $K\times V$ memory is carried between chunks.
\end{enumerate}

There are two useful scheduling choices.  Precomputing $\ve{\eta}_j$ and $\ve{W}_j$ exposes the same
short inter-chunk spine as the conventional compact-WY algorithm, but requires a distinct
$C\times K$ factor $\ve{W}_j$ for every value coordinate.  Forming these factors costs
$O(VC^2K)$ work and storing them costs $O(VCK)$ per chunk.  Alternatively, a kernel can retain
the shared $\ve{L}$ and perform the one-right-hand-side solves only after
$\ve{S}^{\mathrm{in}}$ arrives.  This avoids the value-wise $\ve{W}_j$ tensors and costs
\begin{equation}
  O(C^2K+C^2V+CKV)
  \label{eq:value-chunk-direct-cost}
\end{equation}
work per chunk, but moves the triangular solve onto the inter-chunk dependency path.  A practical
kernel can tile the value dimension, reuse each tile of $\ve{L}$, and use matrix products for the
off-diagonal block updates.
The reverse pass uses the corresponding value-batched unit-upper-triangular solves and has the
same asymptotic work; chunk-local factors may be recomputed rather than saved.

The matrices $\ve{L}$ and $\ve{L}_{qp}$ are shared, but the actual WY operator is not:
\begin{equation}
  \ve{A}_j=(\ve{I}_C+\ve{L}\ve{\Lambda}^{(\beta)}_j)^{-1}.
  \label{eq:value-chunk-sharing-limit}
\end{equation}
Thus arbitrary value-dependent gains require $V$ batched triangular systems.  If
$\beta_{r,j}=\beta_r$ for every $j$, all $\ve{\Lambda}^{(\beta)}_j$ coincide and one shared WY
solve with $V$ right-hand sides recovers the standard KDA/DeltaNet implementation.  Sharing a gain
within each
of $G$ value groups similarly requires only $G$ systems; it is exact for the grouped model and an
approximation to the fully value-wise model.

Finally, this appendix parallelizes the value-wise isotropic boundary, for which the gain stage is
itself a collection of scalar M\"obius scans.  Conditional on known $\ve{g}_t$ and
$\ve{\beta}_t$, the KDN-2 mean admits the same two-key derivation after replacing
$\overline{\ve{p}}_t$ by $\ve{\Gamma}_t^{-1}\ve{g}_t$.  The joint KDN-2 covariance projection,
however, still couples the key and value factors from one token to the next, so this conditional
WY construction does not by itself make the full KDN-2 filter scan parallel.  Here ``exact'' means
exact evaluation of the retained value-wise recurrence, including the trace-projected covariance
prediction in Equation~\eqref{eq:value-prediction}; it does not turn that projected family into the
unrestricted dense Kalman filter.

\section{Two-stage chunkwise implementation of two-axis KDN-2}
\label{app:kdn2-chunkwise}

This appendix derives an exact two-stage implementation of the retained two-axis KDN-2 recurrence
in Section~\ref{sec:kdn2}.  The first stage advances the coupled key--value uncertainty factors and
emits the coefficients required by the mean update.  The second stage conditions on those
coefficients and applies the two-key compact-WY construction.  The distinction is essential: the
mean pass is chunkwise, whereas the joint reverse-KL factor pass remains causal and sequential.

Here ``exact'' means equality with the token-by-token retained KDN-2 recurrence under the same
prediction, factor gauge, and reverse-KL projection---not equality with the unrestricted dense
Kalman filter.

\subsection{Stage 1: causal coupled-factor prepass}
\label{app:kdn2-factor-pass}

The covariance dynamics do not depend on the realized values or on the memory mean.  They may
therefore be evaluated before the $K\times V$ mean recurrence.  Starting from the projected
factors $(\ve{P}_{t-1},\ve{B}_{t-1})$, compute
\begin{equation}
\begin{aligned}
  \widehat{\ve{P}}_t
  &=\ve{D}_t\ve{P}_{t-1}\ve{D}_t^{\tr}
    +\ve{\Omega}^{(k)}_t,
  &
  \widehat{\ve{B}}_t
  &=\ve{B}_{t-1}+\ve{\Omega}^{(v)}_t.
\end{aligned}
  \label{eq:kdn2-chunk-factor-prediction}
\end{equation}
Fix the predictive factor gauge, for example by enforcing
$\det(\widehat{\ve{B}}_t)=1$ and applying the compensating inverse scale to
$\widehat{\ve{P}}_t$.  The exclusive predictive factors emit
\begin{equation}
\begin{aligned}
  \ve{g}_t
  &=\widehat{\ve{P}}_t\ve{k}_t,
  &
  z_t
  &=\ve{k}_t^{\tr}\ve{g}_t,
  \\
  \Delta_{t,j}
  &=r_{t,j}+\widehat b_{t,j}z_t,
  &
  \beta_{t,j}
  &=\frac{\widehat b_{t,j}}{\Delta_{t,j}}.
\end{aligned}
  \label{eq:kdn2-chunk-factor-emissions}
\end{equation}
These quantities are emitted before projecting the current posterior: the current mean update
uses the exclusive factors $\widehat{\ve{P}}_t$ and $\widehat{\ve{B}}_t$, not the post-write
factors.

After this emission, initialize the candidate posterior at the predictive factors and solve the
coupled reverse-KL projection.  One alternating sweep computes
\begin{equation}
\begin{aligned}
  A_P&=\sum_{i=1}^{K}\frac{p_{t,i}}{\widehat p_{t,i}},
  &
  Z_P&=\sum_{i=1}^{K}p_{t,i}k_{t,i}^{2},
  &
  b_{t,j}
  &\leftarrow
  \frac{K}{A_P/\widehat b_{t,j}+Z_P/r_{t,j}},
  \\
  A_B&=\sum_{j=1}^{V}\frac{b_{t,j}}{\widehat b_{t,j}},
  &
  T_B&=\sum_{j=1}^{V}\frac{b_{t,j}}{r_{t,j}},
  &
  p_{t,i}
  &\leftarrow
  \frac{V}{A_B/\widehat p_{t,i}+T_Bk_{t,i}^{2}}.
\end{aligned}
  \label{eq:kdn2-chunk-factor-projection}
\end{equation}
Gauge normalization follows each complete sweep.  Iteration to convergence gives the unique
reverse-KL optimum in the fixed gauge.  Alternatively, a fixed number
$N_{\mathrm{proj}}$ of sweeps defines a truncated-projection KDN-2 recurrence.  In either case,
the resulting $(\ve{P}_t,\ve{B}_t)$ are carried to token $t+1$, while
$\ve{g}_t$, $\ve{\beta}_t$, and optionally $\ve{\Delta}_t$ are stored for Stage~2.

\paragraph{Exact scalarization of the coupled projection.}
The apparent $(K+V)$-dimensional fixed point has a one-dimensional representation because the
balancing matrix in Equation~\eqref{eq:kdn2-balancing-matrix} is the sum of two rank-one terms.
Define the predictive axis statistics and factor-summary ratios
\begin{equation}
  \zeta^{(k)}_{t,i}=\widehat p_{t,i}k_{t,i}^{2},
  \qquad
  \zeta^{(v)}_{t,j}=\frac{\widehat b_{t,j}}{r_{t,j}},
  \qquad
  \tau^{(k)}_t=\frac{Z_P}{A_P},
  \qquad
  \tau^{(v)}_t=\frac{T_B}{A_B}.
  \label{eq:kdn2-chunk-scalar-statistics}
\end{equation}
Let $\ve{\zeta}^{(k)}_t\in\R_+^K$ and
$\ve{\zeta}^{(v)}_t\in\R_+^V$ collect the corresponding axis statistics.
For a nonnegative vector $\ve{a}\in\R_+^m$, introduce
\begin{equation}
\begin{aligned}
  D_{\ve{a}}(s)&=\sum_{\ell=1}^{m}\frac{1}{1+s a_\ell},
  &
  N_{\ve{a}}(s)&=\sum_{\ell=1}^{m}\frac{a_\ell}{1+s a_\ell},
  &
  F_{\ve{a}}(s)&=\frac{N_{\ve{a}}(s)}{D_{\ve{a}}(s)}.
\end{aligned}
  \label{eq:kdn2-chunk-weighted-map}
\end{equation}
The stationary equations in Equation~\eqref{eq:kdn2-chunk-factor-projection} imply
\begin{equation}
  \frac{b_{t,j}}{\widehat b_{t,j}}
  =\frac{K/A_P}{1+\tau^{(k)}_t\zeta^{(v)}_{t,j}},
  \qquad
  \frac{p_{t,i}}{\widehat p_{t,i}}
  =\frac{V/A_B}{1+\tau^{(v)}_t\zeta^{(k)}_{t,i}}.
  \label{eq:kdn2-chunk-factor-ratio-form}
\end{equation}
Taking the defining ratios reduces the joint projection to
\begin{equation}
  \tau^{(v)}_t=F_{\ve{\zeta}^{(v)}_t}(\tau^{(k)}_t),
  \qquad
  \tau^{(k)}_t=F_{\ve{\zeta}^{(k)}_t}(\tau^{(v)}_t),
  \qquad\Longleftrightarrow\qquad
  \tau^{(k)}_t=F_{\ve{\zeta}^{(k)}_t}\!\left(
    F_{\ve{\zeta}^{(v)}_t}(\tau^{(k)}_t)
  \right).
  \label{eq:kdn2-chunk-scalar-root}
\end{equation}
Thus the converged reverse-KL projection requires one scalar root rather than alternating writes
of all $K+V$ candidate factors.  Define
$G_t(s)=F_{\ve{\zeta}^{(k)}_t}(F_{\ve{\zeta}^{(v)}_t}(s))$ and
$\mathcal{R}_t(s)=G_t(s)-s$.  Since $G_t$ is a weighted mean of the entries of
$\ve{\zeta}^{(k)}_t$,
\begin{equation}
  \mathcal{R}_t\!\left(\min_i\zeta^{(k)}_{t,i}\right)\ge0,
  \qquad
  \mathcal{R}_t\!\left(\max_i\zeta^{(k)}_{t,i}\right)\le0.
  \label{eq:kdn2-chunk-scalar-bracket}
\end{equation}
Uniqueness follows from the strict convexity of the gauge-fixed reverse-KL objective.  If all
$\zeta^{(k)}_{t,i}$ equal one constant, that constant is the root directly; otherwise safeguarded
Newton or bisection applies on the bracket above.  For deterministic GPU execution, use a fixed
maximum iteration count and the gauge-invariant relative residual
$|G_t(s)-s|/(G_t(s)+s)\le\epsilon_{\mathrm{root}}$, with the all-zero case handled directly.

After solving for $\tau^{(k)}_t$ and setting
$\tau^{(v)}_t=F_{\ve{\zeta}^{(v)}_t}(\tau^{(k)}_t)$, one convenient raw gauge is
\begin{equation}
  p_{t,i}^{(0)}
  =\frac{\widehat p_{t,i}}{1+\tau^{(v)}_t\zeta^{(k)}_{t,i}},
  \qquad
  b_{t,j}^{(0)}
  =\frac{(1+\tau^{(k)}_t\tau^{(v)}_t)\widehat b_{t,j}}
  {1+\tau^{(k)}_t\zeta^{(v)}_{t,j}}.
  \label{eq:kdn2-chunk-scalar-reconstruction}
\end{equation}
Every gauge-equivalent solution is
$p_{t,i}=p_{t,i}^{(0)}/\lambda_t$ and
$b_{t,j}=\lambda_t b_{t,j}^{(0)}$.  For the determinant gauge in
Equation~\eqref{eq:kdn2-gauge-fix}, choose
\begin{equation}
  \lambda_t
  =\exp\!\left(-\frac{1}{V}\sum_{j=1}^{V}\log b_{t,j}^{(0)}\right).
  \label{eq:kdn2-chunk-scalar-gauge}
\end{equation}
The gauge-free entry variances are
\begin{equation}
  p_{t,i}b_{t,j}
  =\frac{(1+\tau^{(k)}_t\tau^{(v)}_t)\widehat p_{t,i}\widehat b_{t,j}}
  {(1+\tau^{(v)}_t\zeta^{(k)}_{t,i})
   (1+\tau^{(k)}_t\zeta^{(v)}_{t,j})}.
  \label{eq:kdn2-chunk-scalar-entry-variance}
\end{equation}
The identity $D_{\ve{a}}(s)+sN_{\ve{a}}(s)=m$ verifies the fixed-point form of
Equation~\eqref{eq:kdn2-chunk-factor-projection} by direct substitution.

The scalar form also preserves a deliberately truncated projection.  Predictive initialization
gives
\begin{equation}
  \tau_t^{(k,0)}=\frac{1}{K}\sum_{i=1}^{K}\zeta^{(k)}_{t,i},
  \qquad
  \tau_t^{(v,n)}=F_{\ve{\zeta}^{(v)}_t}(\tau_t^{(k,n)}),
  \qquad
  \tau_t^{(k,n+1)}=F_{\ve{\zeta}^{(k)}_t}(\tau_t^{(v,n)}).
  \label{eq:kdn2-chunk-truncated-scalar-iteration}
\end{equation}
Here $\tau_t^{(k,n)}$ is the incoming $p$-summary used by the $b$ update in sweep $n$, whereas
$\tau_t^{(k,n+1)}$ is the outgoing $p$-summary.  For $N_{\mathrm{proj}}\ge1$, the last outgoing
summary need not be evaluated: use the lagged pair
$(\tau_t^{(k,N_{\mathrm{proj}}-1)},\tau_t^{(v,N_{\mathrm{proj}}-1)})$ in
Equation~\eqref{eq:kdn2-chunk-scalar-reconstruction}.  Indeed,
$D_{\ve{\zeta}^{(v)}_t}(s)+sN_{\ve{\zeta}^{(v)}_t}(s)=V$ and
$\tau^{(v)}=N_{\ve{\zeta}^{(v)}_t}(s)/D_{\ve{\zeta}^{(v)}_t}(s)$ imply
$V/D_{\ve{\zeta}^{(v)}_t}(s)=1+s\tau^{(v)}$, precisely the reconstruction prefactor.  The result
is therefore the same posterior covariance, up to gauge, as $N_{\mathrm{proj}}$ complete
$b$-then-$p$ sweeps, not a new approximation.  The case $N_{\mathrm{proj}}=0$ simply returns the
predictive factors.

On a GPU, $\ve{\zeta}^{(k)}_t$ and $\ve{\zeta}^{(v)}_t$ can remain on chip while each scalar
iteration performs one reduction along each axis.  The vectors $\ve{p}_t$ and $\ve{b}_t$ are
reconstructed and gauge-normalized only once per token.  A persistent block for each independent
batch--head pair can fuse prediction, gain emission, scalar projection, reconstruction, and gauge
normalization.  For a converged root, implicit differentiation can avoid saving or unrolling the
root iterations; a fixed truncated solver should instead be differentiated through to preserve
its specified recurrence.  These changes can reduce synchronization and do reduce factor-vector
traffic inside a token, but they do not remove the causal dependence between tokens.

\paragraph{Why Stage 1 is not an associative scan.}
The map
\begin{equation}
  (\ve{P}_{t-1},\ve{B}_{t-1})
  \longmapsto
  (\ve{P}_t,\ve{B}_t)
  \label{eq:kdn2-chunk-factor-map}
\end{equation}
contains a joint nonlinear matrix-scaling solve whose summaries depend on both candidate factors.
Unlike the $2\times2$ M\"obius matrices of the one-axis recurrence, composing token maps does not
produce another token-local map in the same $O(K+V)$-parameter family.  Stage~1 chunks must
therefore be traversed in causal order; chunking provides tiling and checkpointing, not time
parallelism.

\subsection{Stage 2: conditional two-key mean recurrence}
\label{app:kdn2-conditional-mean}

Once Stage~1 has produced $\ve{g}_t$ and $\ve{\beta}_t$, the memory recurrence is
\begin{equation}
  \ve{S}_t
  =\ve{D}_t\ve{S}_{t-1}
   +\ve{g}_t
    \left[
      \ve{\beta}_t\odot
      \left(
        \ve{v}_t-(\ve{D}_t\ve{S}_{t-1})^{\tr}\ve{k}_t
      \right)
    \right]^{\tr}.
  \label{eq:kdn2-chunk-original-mean}
\end{equation}
Consider one chunk of length $C$, relabel its tokens by $r=1,\ldots,C$, and let
$\ve{S}^{\mathrm{in}}=\ve{S}_0$ denote its incoming physical memory.  Assume positive diagonal
decays and define
\begin{equation}
  \ve{\Gamma}_r
  =\ve{D}_r\ve{D}_{r-1}\cdots\ve{D}_1,
  \qquad
  \ve{\Gamma}_0=\ve{I}_K,
  \qquad
  \ve{S}_r=\ve{\Gamma}_r\ve{Y}_r.
  \label{eq:kdn2-chunk-cumulative-decay}
\end{equation}
The decay-scaled read and base-write directions are
\begin{equation}
  \ve{k}^{+}_r=\ve{\Gamma}_r^{\tr}\ve{k}_r,
  \qquad
  \ve{g}^{-}_r=\ve{\Gamma}_r^{-1}\ve{g}_r.
  \label{eq:kdn2-chunk-two-keys}
\end{equation}
Since $\ve{D}_r\ve{\Gamma}_{r-1}=\ve{\Gamma}_r$, multiplying
Equation~\eqref{eq:kdn2-chunk-original-mean} by $\ve{\Gamma}_r^{-1}$ gives
\begin{equation}
  \ve{Y}_r
  =\ve{Y}_{r-1}
   +\ve{g}^{-}_r
    \left[
      \ve{\beta}_r\odot
      \left(
        \ve{v}_r-\ve{Y}_{r-1}^{\tr}\ve{k}^{+}_r
      \right)
    \right]^{\tr},
  \qquad
  \ve{Y}_0=\ve{S}^{\mathrm{in}}.
  \label{eq:kdn2-chunk-normalized-mean}
\end{equation}
Thus normalized KDN-2 is a two-key delta rule: $\ve{k}^{+}_r$ reads, while
$\beta_{r,j}\ve{g}^{-}_r$ writes value coordinate $j$.

\subsection{Value-batched causal systems and compact-WY form}
\label{app:kdn2-triangular-systems}

Stack the scaled directions by row,
\begin{equation}
  [\ve{K}_{+}]_{r,:}=(\ve{k}^{+}_r)^{\tr},
  \qquad
  [\ve{G}_{-}]_{r,:}=(\ve{g}^{-}_r)^{\tr},
  \qquad
  \ve{L}
  =\operatorname{tril}(\ve{K}_{+}\ve{G}_{-}^{\tr},-1).
  \label{eq:kdn2-chunk-shared-interaction}
\end{equation}
The strictly lower interaction matrix $\ve{L}\in\R^{C\times C}$ is shared across value
coordinates.  For column $j$, define
\begin{equation}
\begin{aligned}
  \ve{v}_j
  &=(v_{1,j},\ldots,v_{C,j})^{\tr},
  &
  \ve{\Lambda}_j
  &=\diag(\beta_{1,j},\ldots,\beta_{C,j}),
  \\
  \varepsilon_{r,j}
  &=v_{r,j}-(\ve{k}^{+}_r)^{\tr}\ve{y}_{r-1,j},
  &
  \ve{\xi}_j
  &=\ve{\Lambda}_j\ve{\varepsilon}_j.
\end{aligned}
  \label{eq:kdn2-chunk-innovation-variables}
\end{equation}
Unrolling only writes from earlier positions gives
\begin{equation}
  \ve{\varepsilon}_j
  =\ve{v}_j-\ve{K}_{+}\ve{s}^{\mathrm{in}}_j-\ve{L}\ve{\xi}_j,
  \qquad
  \ve{\xi}_j=\ve{\Lambda}_j\ve{\varepsilon}_j.
  \label{eq:kdn2-chunk-causal-unrolling}
\end{equation}
Therefore the causal innovations are the unique solutions of either equivalent unit-lower system
\begin{equation}
\begin{aligned}
  (\ve{I}_C+\ve{L}\ve{\Lambda}_j)\ve{\varepsilon}_j
  &=\ve{v}_j-\ve{K}_{+}\ve{s}^{\mathrm{in}}_j,
  \\
  (\ve{I}_C+\ve{\Lambda}_j\ve{L})\ve{\xi}_j
  &=\ve{\Lambda}_j
    (\ve{v}_j-\ve{K}_{+}\ve{s}^{\mathrm{in}}_j).
\end{aligned}
  \label{eq:kdn2-chunk-triangular-systems}
\end{equation}
The first form remains well defined when some gains are zero.  Let
$\ve{\Xi}=[\ve{\xi}_1,\ldots,\ve{\xi}_V]\in\R^{C\times V}$ and let
$\ve{\xi}_r^{\tr}$ denote its $r$th row.  The value-batched forward substitution is
\begin{equation}
  \ve{\xi}_r
  =\ve{\beta}_r\odot
   \left[
     \ve{v}_r
     -(\ve{S}^{\mathrm{in}})^{\tr}\ve{k}^{+}_r
     -\sum_{i<r}L_{ri}\ve{\xi}_i
   \right].
  \label{eq:kdn2-chunk-fused-substitution}
\end{equation}
The geometry $\ve{L}$ is shared, but its effective columns are scaled differently for every value
coordinate.

To expose the compact-WY factors, define
\begin{equation}
  \ve{A}_j
  =(\ve{I}_C+\ve{L}\ve{\Lambda}_j)^{-1},
  \qquad
  \ve{\eta}_j=\ve{A}_j\ve{v}_j,
  \qquad
  \ve{W}_j=\ve{A}_j\ve{K}_{+}.
  \label{eq:kdn2-chunk-wy-factors}
\end{equation}
The inverse denotes a unit-lower-triangular solve.  For the normalized rank-one transition
\begin{equation}
  \ve{E}_{r,j}
  =\ve{I}_K
   -\beta_{r,j}\ve{g}^{-}_r(\ve{k}^{+}_r)^{\tr},
  \label{eq:kdn2-chunk-rank-one-transition}
\end{equation}
the compact-WY identity is
\begin{equation}
  \ve{E}_{C,j}\ve{E}_{C-1,j}\cdots\ve{E}_{1,j}
  =\ve{I}_K-\ve{G}_{-}^{\tr}\ve{\Lambda}_j\ve{W}_j.
  \label{eq:kdn2-chunk-compact-wy-identity}
\end{equation}
The unscaled innovations satisfy
\begin{equation}
  \ve{\varepsilon}_j
  =\ve{\eta}_j-\ve{W}_j\ve{s}^{\mathrm{in}}_j,
  \qquad
  \ve{\xi}_j
  =\ve{\Lambda}_j
   (\ve{\eta}_j-\ve{W}_j\ve{s}^{\mathrm{in}}_j).
  \label{eq:kdn2-chunk-wy-innovation}
\end{equation}

The innovations $\varepsilon_{r,j}$ are exactly the exclusive residuals $\delta_{t,j}$ in
Equation~\eqref{eq:kdn2-key-statistics}.  Hence the predictive likelihood can be evaluated from
$\ve{\varepsilon}$ and the variances $\ve{\Delta}$ emitted by Stage~1 without another recurrent
read.

\subsection{Parallel reads, chunk carry, and stable ratios}
\label{app:kdn2-chunk-reads}

For a post-write query $\ve{q}_r$, define
\begin{equation}
  [\ve{Q}_{+}]_{r,:}
  =(\ve{\Gamma}_r^{\tr}\ve{q}_r)^{\tr},
  \qquad
  \ve{L}_{qg}
  =\operatorname{tril}(\ve{Q}_{+}\ve{G}_{-}^{\tr}).
  \label{eq:kdn2-chunk-query-interaction}
\end{equation}
The diagonal is retained because the current write is visible to the current query.  All outputs
and the physical outgoing state are
\begin{equation}
\begin{aligned}
  \ve{O}
  &=\ve{Q}_{+}\ve{S}^{\mathrm{in}}+\ve{L}_{qg}\ve{\Xi},
  \\
  \ve{S}^{\mathrm{out}}
  &=\ve{\Gamma}_C
    \left(
      \ve{S}^{\mathrm{in}}+\ve{G}_{-}^{\tr}\ve{\Xi}
    \right).
\end{aligned}
  \label{eq:kdn2-chunk-outputs-state}
\end{equation}
A pre-write read replaces $\operatorname{tril}$ by its strictly lower part.

The inverse cumulative decay in $\ve{g}^{-}_r$ need not be materialized.  Define the physical
transition from immediately after token $i$ to token $r$ by
\begin{equation}
  \ve{\Phi}_{r\leftarrow i}
  =
  \begin{cases}
    \ve{D}_r\ve{D}_{r-1}\cdots\ve{D}_{i+1}, & r>i,\\
    \ve{I}_K, & r=i.
  \end{cases}
  \label{eq:kdn2-chunk-decay-ratio}
\end{equation}
Then all required interactions have the stable form
\begin{equation}
\begin{aligned}
  L_{ri}
  &=\ve{k}_r^{\tr}\ve{\Phi}_{r\leftarrow i}\ve{g}_i,
  && i<r,
  \\
  (L_{qg})_{ri}
  &=\ve{q}_r^{\tr}\ve{\Phi}_{r\leftarrow i}\ve{g}_i,
  && i\le r,
  \\
  \ve{S}^{\mathrm{out}}
  &=\ve{\Gamma}_C\ve{S}^{\mathrm{in}}
    +\sum_{i=1}^{C}
      \ve{\Phi}_{C\leftarrow i}\ve{g}_i\ve{\xi}_i^{\tr}.
\end{aligned}
  \label{eq:kdn2-chunk-stable-ratios}
\end{equation}
For positive diagonal decays, $\ve{\Phi}_{r\leftarrow i}$ should be formed from differences of
cumulative log decays.  The decay applied to a write at token $i$ begins at token $i+1$;
$\ve{D}_i$ acts before that write.

Under a time-local gauge transformation
\begin{equation}
  \widehat{\ve{B}}_i\mapsto\lambda_i\widehat{\ve{B}}_i,
  \qquad
  \widehat{\ve{P}}_i\mapsto\widehat{\ve{P}}_i/\lambda_i,
  \label{eq:kdn2-chunk-local-gauge}
\end{equation}
we have $\ve{g}_i\mapsto\ve{g}_i/\lambda_i$ and
$\ve{\beta}_i\mapsto\lambda_i\ve{\beta}_i$.  Hence
$L_{ri}\beta_{i,j}$, $\ve{g}_i\xi_{i,j}$, and all physical outputs in
Equations~\eqref{eq:kdn2-chunk-triangular-systems}--\eqref{eq:kdn2-chunk-stable-ratios} are gauge
invariant.

\subsection{Two-stage schedule, cost, and exactness}
\label{app:kdn2-chunk-schedule}

The complete forward schedule is:
\begin{enumerate}
  \item \textbf{Coupled-factor stage.} Traverse tokens causally with
    Equations~\eqref{eq:kdn2-chunk-factor-prediction}--\eqref{eq:kdn2-chunk-factor-emissions},
    and evaluate the projection either by alternating
    Equation~\eqref{eq:kdn2-chunk-factor-projection} or by the scalar form in
    Equations~\eqref{eq:kdn2-chunk-scalar-root}--\eqref{eq:kdn2-chunk-truncated-scalar-iteration}.
    Emit $\ve{g}_t$ and $\ve{\beta}_t$, optionally retain $\ve{\Delta}_t$ for the predictive
    likelihood, and carry the projected factors.
  \item \textbf{Conditional-mean stage.} For each memory chunk, form the two-key interactions in
    Equation~\eqref{eq:kdn2-chunk-shared-interaction}, solve
    Equation~\eqref{eq:kdn2-chunk-triangular-systems} batched over value coordinates, evaluate
    Equation~\eqref{eq:kdn2-chunk-outputs-state}, and carry $\ve{S}^{\mathrm{out}}$ to the next
    chunk.
\end{enumerate}

With $N_{\mathrm{proj}}$ reverse-KL sweeps per token, Stage~1 costs
\begin{equation}
  O\!\left(T(N_{\mathrm{proj}}+1)(K+V)\right)
  \label{eq:kdn2-chunk-factor-cost}
\end{equation}
work and has sequential depth proportional to $T(N_{\mathrm{proj}}+1)$.  Storing the factored gain
sequence costs $O(T(K+V))$; checkpointing may instead recompute it during backpropagation.
The scalar evaluation has the same asymptotic reduction work, with $N_{\mathrm{proj}}$ scalar
sweep steps but only $N_{\mathrm{proj}}-1$ required key-ratio updates, and reconstructs the factor
vectors once per token.  For a converged solve, replace $N_{\mathrm{proj}}$ by the number
$N_{\mathrm{root}}$ of safeguarded scalar root iterations.  In either case the token-to-token
depth remains linear in $T$.

For one Stage~2 chunk, constructing the shared interactions costs $O(C^2K)$, the value-batched
triangular systems cost $O(C^2V)$, and multiplication by the carried memory costs $O(CKV)$.
Across the sequence, the direct implementation costs
\begin{equation}
  O\!\left(TC(K+V)+TKV\right)
  \label{eq:kdn2-chunk-mean-cost}
\end{equation}
with $O(KV+C^2+CV)$ working storage beyond inputs and outputs.  In the direct schedule, each
chunk's $C$-row solve waits for $\ve{S}^{\mathrm{in}}$, so the per-trajectory critical path remains
$\Theta(T)$, vectorized over values.  As in Appendix~\ref{app:value-wise-chunkwise}, precomputing
distinct $\ve{W}_j$ factors moves those local solves off the critical path and leaves $T/C$ chunk
carries, but costs $O(VC^2K)$ work and $O(VCK)$ storage per chunk.  If
$\beta_{r,j}=\beta_r$ for every $j$, Stage~2 reduces to one shared DeltaNet/KDA triangular system
with $V$ right-hand sides.  Tying gains within $G$ value groups requires $G$ systems and is exact
only for the corresponding grouped uncertainty model.

\paragraph{Exactness.}
Assume positive diagonal decays and let Stage~1 use the same prediction, gauge, and reverse-KL
solution as the sequential KDN-2 reference.  Covariance updates are independent of the realized
observation, so Stage~1 may run ahead of the mean recurrence and emits the same $\ve{g}_t$ and
$\ve{\beta}_t$ by causal induction.  For fixed gains,
Equation~\eqref{eq:kdn2-chunk-normalized-mean} gives
$\ve{y}_{r-1,j}=\ve{s}^{\mathrm{in}}_j+
\sum_{i<r}\ve{g}^{-}_i\xi_{i,j}$.  Substitution into $\varepsilon_{r,j}$ yields the unit-lower
system in Equation~\eqref{eq:kdn2-chunk-causal-unrolling}, whose unique solution equals the
sequential innovations.  Substituting back gives the same prefix and outgoing states; applying
the queries gives Equation~\eqref{eq:kdn2-chunk-outputs-state}.  Induction over chunk boundaries
therefore proves equality of posterior means, predictive residuals, outputs, and boundary states.

\end{document}
